% ============================================================
% Section 118: 超导环持久电流与电阻率上限
% ============================================================
\section{固体物理：超导环持久电流与电阻率上限}
\label{sec:persistent-current}

\begin{modelbox}[题目：持久电流实验估计超导电阻率]
一个环由直径 $d=1\,\text{mm}$ 的铅丝弯成直径 $D=10\,\text{cm}$ 的圆。环处于超导态并通有 $100\,\text{A}$ 的电流，一年内未观测到可探测的变化。若检测器对 $1\,\mu\text{A}$ 的变化都能检测到，计算实验上超导态铅的电阻率上限。（答案近似因子不超过 3 即可。）
\end{modelbox}

\begin{tipbox}[考点快速分析]
\begin{itemize}
    \item \textbf{RL 衰减}：$I(t)=I(0)e^{-Rt/L}$，微小衰减 $\Delta I/I(0)\ll1$ 时对数近似
    \item \textbf{自感估算}：细圆环近似 $L\approx(\pi/2)\mu_0r$
    \item \textbf{电阻率}：$\rho=RS/l$，$S$ 为截面积，$l$ 为丝长
    \item \textbf{核心思想}：用电流不衰减的时间尺度反推电阻率的极低上限
\end{itemize}
\end{tipbox}

\begin{derivationbox}[标准解题过程]
\textbf{1. RL 电路衰减模型}

将超导环视为电感 $L$、电阻 $R$ 的 RL 串联电路。无电源时电流指数衰减：
\begin{equation}
I(t)=I(0)e^{-Rt/L}.
\end{equation}

检测器可分辨的最小电流变化为 $\Delta I=1\,\mu\text{A}$，一年内未观测到变化，说明 $\Delta I\le1\,\mu\text{A}$。因此
\begin{equation}
\frac{I(0)}{I(t)}=\frac{I(0)}{I(0)-\Delta I}\approx1+\frac{\Delta I}{I(0)}=1+10^{-8},
\end{equation}
\begin{equation}
\ln\!\frac{I(0)}{I(t)}\approx\frac{\Delta I}{I(0)}=10^{-8}.
\end{equation}

由 $R=(L/t)\ln[I(0)/I(t)]$ 得
\begin{equation}
R\approx\frac{L}{t}\cdot10^{-8}.
\end{equation}

\textbf{2. 电感 $L$ 的近似}

铅丝环半径 $r=D/2=0.05\,\text{m}$。细圆环近似下，环内磁场均匀且等于环中心磁场：
\begin{equation}
B\approx\frac{\mu_0I}{2r}.
\end{equation}

环面积 $S_2=\pi r^2$，磁通量
\begin{equation}
\Phi\approx BS_2=\frac{\mu_0I}{2r}\cdot\pi r^2=\frac{\pi}{2}\mu_0Ir.
\end{equation}

自感系数
\begin{equation}
L=\frac{\Phi}{I}=\frac{\pi}{2}\mu_0r.
\end{equation}

代入 $\mu_0=4\pi\times10^{-7}\,\text{H/m}$，$r=0.05\,\text{m}$：
\begin{equation}
L=\frac{\pi}{2}\times4\pi\times10^{-7}\times0.05=\pi^2\times10^{-8}\approx9.87\times10^{-8}\,\text{H}.
\end{equation}

\textbf{3. 电阻 $R$ 的上限}

时间 $t=1\,\text{年}\approx3.15\times10^7\,\text{s}$。故
\begin{equation}
R=\frac{9.87\times10^{-8}}{3.15\times10^7}\times10^{-8}\approx3.13\times10^{-23}\,\Omega.
\end{equation}

\textbf{4. 电阻率 $\rho$ 的上限}

铅丝长度 $l=2\pi r\approx0.314\,\text{m}$，横截面积 $S_1=\pi(d/2)^2=\pi\times(0.5\times10^{-3})^2\approx7.85\times10^{-7}\,\text{m}^2$。

由 $R=\rho l/S_1$：
\begin{equation}
\rho=R\frac{S_1}{l}=3.13\times10^{-23}\times\frac{7.85\times10^{-7}}{0.314}\approx7.8\times10^{-29}\,\Omega\cdot\text{m}.
\end{equation}

\begin{equation}
\boxed{\rho\le7.8\times10^{-29}\,\Omega\cdot\text{m}.}
\end{equation}
\end{derivationbox}

\begin{formulabox}[答案汇总]
\begin{center}
\begin{tabular}{ll}
\toprule
\textbf{项目} & \textbf{结果} \\
\midrule
自感（细圆环近似） & $L\approx(\pi/2)\mu_0r\approx9.9\times10^{-8}\,\text{H}$ \\
电阻上限 & $R\lesssim3.1\times10^{-23}\,\Omega$ \\
电阻率上限 & $\rho\le7.8\times10^{-29}\,\Omega\cdot\text{m}$ \\
对比（室温 Cu） & $\rho_{\text{Cu}}\approx1.7\times10^{-8}\,\Omega\cdot\text{m}$，超导 Pb 低 21 个量级 \\
\bottomrule
\end{tabular}
\end{center}
\end{formulabox}

\begin{insightbox}[物理图像]
\begin{itemize}
    \item 持久电流实验是验证超导体零电阻最直接、最有力的证据。通过观测电流不衰减的时间尺度，可给出电阻率的极低上限。
    \item 超导态铅的电阻率上限比室温铜低约 21 个数量级。这意味着即便电阻率不为零，在宇宙年龄尺度上也无法观测到电流衰减。
    \item 电感计算采用细圆环近似，在允许误差因子内足够精确。更严格的计算需考虑导线截面的电流分布，但数量级不变。
\end{itemize}
\end{insightbox}
